% Options for packages loaded elsewhere
\PassOptionsToPackage{unicode}{hyperref}
\PassOptionsToPackage{hyphens}{url}
\documentclass[
]{article}
\usepackage{xcolor}
\usepackage[margin=1in]{geometry}
\usepackage{amsmath,amssymb}
\setcounter{secnumdepth}{-\maxdimen} % remove section numbering
\usepackage{iftex}
\ifPDFTeX
  \usepackage[T1]{fontenc}
  \usepackage[utf8]{inputenc}
  \usepackage{textcomp} % provide euro and other symbols
\else % if luatex or xetex
  \usepackage{unicode-math} % this also loads fontspec
  \defaultfontfeatures{Scale=MatchLowercase}
  \defaultfontfeatures[\rmfamily]{Ligatures=TeX,Scale=1}
\fi
\usepackage{lmodern}
\ifPDFTeX\else
  % xetex/luatex font selection
\fi
% Use upquote if available, for straight quotes in verbatim environments
\IfFileExists{upquote.sty}{\usepackage{upquote}}{}
\IfFileExists{microtype.sty}{% use microtype if available
  \usepackage[]{microtype}
  \UseMicrotypeSet[protrusion]{basicmath} % disable protrusion for tt fonts
}{}
\makeatletter
\@ifundefined{KOMAClassName}{% if non-KOMA class
  \IfFileExists{parskip.sty}{%
    \usepackage{parskip}
  }{% else
    \setlength{\parindent}{0pt}
    \setlength{\parskip}{6pt plus 2pt minus 1pt}}
}{% if KOMA class
  \KOMAoptions{parskip=half}}
\makeatother
\usepackage{color}
\usepackage{fancyvrb}
\newcommand{\VerbBar}{|}
\newcommand{\VERB}{\Verb[commandchars=\\\{\}]}
\DefineVerbatimEnvironment{Highlighting}{Verbatim}{commandchars=\\\{\}}
% Add ',fontsize=\small' for more characters per line
\newenvironment{Shaded}{}{}
\newcommand{\AlertTok}[1]{\textcolor[rgb]{1.00,0.00,0.00}{\textbf{#1}}}
\newcommand{\AnnotationTok}[1]{\textcolor[rgb]{0.38,0.63,0.69}{\textbf{\textit{#1}}}}
\newcommand{\AttributeTok}[1]{\textcolor[rgb]{0.49,0.56,0.16}{#1}}
\newcommand{\BaseNTok}[1]{\textcolor[rgb]{0.25,0.63,0.44}{#1}}
\newcommand{\BuiltInTok}[1]{\textcolor[rgb]{0.00,0.50,0.00}{#1}}
\newcommand{\CharTok}[1]{\textcolor[rgb]{0.25,0.44,0.63}{#1}}
\newcommand{\CommentTok}[1]{\textcolor[rgb]{0.38,0.63,0.69}{\textit{#1}}}
\newcommand{\CommentVarTok}[1]{\textcolor[rgb]{0.38,0.63,0.69}{\textbf{\textit{#1}}}}
\newcommand{\ConstantTok}[1]{\textcolor[rgb]{0.53,0.00,0.00}{#1}}
\newcommand{\ControlFlowTok}[1]{\textcolor[rgb]{0.00,0.44,0.13}{\textbf{#1}}}
\newcommand{\DataTypeTok}[1]{\textcolor[rgb]{0.56,0.13,0.00}{#1}}
\newcommand{\DecValTok}[1]{\textcolor[rgb]{0.25,0.63,0.44}{#1}}
\newcommand{\DocumentationTok}[1]{\textcolor[rgb]{0.73,0.13,0.13}{\textit{#1}}}
\newcommand{\ErrorTok}[1]{\textcolor[rgb]{1.00,0.00,0.00}{\textbf{#1}}}
\newcommand{\ExtensionTok}[1]{#1}
\newcommand{\FloatTok}[1]{\textcolor[rgb]{0.25,0.63,0.44}{#1}}
\newcommand{\FunctionTok}[1]{\textcolor[rgb]{0.02,0.16,0.49}{#1}}
\newcommand{\ImportTok}[1]{\textcolor[rgb]{0.00,0.50,0.00}{\textbf{#1}}}
\newcommand{\InformationTok}[1]{\textcolor[rgb]{0.38,0.63,0.69}{\textbf{\textit{#1}}}}
\newcommand{\KeywordTok}[1]{\textcolor[rgb]{0.00,0.44,0.13}{\textbf{#1}}}
\newcommand{\NormalTok}[1]{#1}
\newcommand{\OperatorTok}[1]{\textcolor[rgb]{0.40,0.40,0.40}{#1}}
\newcommand{\OtherTok}[1]{\textcolor[rgb]{0.00,0.44,0.13}{#1}}
\newcommand{\PreprocessorTok}[1]{\textcolor[rgb]{0.74,0.48,0.00}{#1}}
\newcommand{\RegionMarkerTok}[1]{#1}
\newcommand{\SpecialCharTok}[1]{\textcolor[rgb]{0.25,0.44,0.63}{#1}}
\newcommand{\SpecialStringTok}[1]{\textcolor[rgb]{0.73,0.40,0.53}{#1}}
\newcommand{\StringTok}[1]{\textcolor[rgb]{0.25,0.44,0.63}{#1}}
\newcommand{\VariableTok}[1]{\textcolor[rgb]{0.10,0.09,0.49}{#1}}
\newcommand{\VerbatimStringTok}[1]{\textcolor[rgb]{0.25,0.44,0.63}{#1}}
\newcommand{\WarningTok}[1]{\textcolor[rgb]{0.38,0.63,0.69}{\textbf{\textit{#1}}}}
\setlength{\emergencystretch}{3em} % prevent overfull lines
\providecommand{\tightlist}{%
  \setlength{\itemsep}{0pt}\setlength{\parskip}{0pt}}
% Preamble for the pdfLaTeX build of schur-six-538-note-draft.md
% (scripts/build-paper-538.sh). The Lean block uses six Unicode symbols;
% pdfLaTeX needs a definition for each.
\usepackage{newunicodechar}
\newunicodechar{→}{\ensuremath{\rightarrow}}
\newunicodechar{≤}{\ensuremath{\le}}
\newunicodechar{∃}{\ensuremath{\exists}}
\newunicodechar{∀}{\ensuremath{\forall}}
\newunicodechar{¬}{\ensuremath{\neg}}
\newunicodechar{∧}{\ensuremath{\wedge}}
\usepackage{bookmark}
\IfFileExists{xurl.sty}{\usepackage{xurl}}{} % add URL line breaks if available
\urlstyle{same}
% fallback for those not using the hyperref driver hyperxmp:
\makeatletter
\@ifundefined{xmpquote}{\newcommand{\xmpquote}[1]{#1}}{}
\makeatother
\hypersetup{
  pdftitle={The sixth Schur number is at least 538},
  pdfauthor={Adam McKenna},
  hidelinks,
  pdfcreator={LaTeX via pandoc}}

\title{The sixth Schur number is at least 538}
\author{Adam McKenna}
\date{November 2026. doi:10.5281/zenodo.23937655}

\begin{document}
\maketitle
\begin{abstract}
We give an explicit partition of \(\{1,\dots,538\}\) into six sum-free
sets. This shows \(S(6)\ge 538\) and improves the bound \(S(6)\ge 537\)
of Chen, Zhu and Long. It also gives \(R_6(3)\ge 540\) for the Ramsey
number of the triangle in six colours, and known recursive constructions
give \(S(8)\ge 5382\), \(S(11)\ge 204588\) and \(S(13)\ge 2045882\). The
partition of \(\{1,\dots,538\}\) is checked directly and in Lean 4.
\end{abstract}

\section{1. The result}\label{the-result}

A set \(A\) of positive integers is \emph{sum-free} when there are no
\(x,y\in A\) with \(x+y\in A\); the elements \(x\) and \(y\) may be
equal. The Schur number \(S(k)\) is the largest \(n\) such that
\(\{1,\dots,n\}\) is the union of \(k\) sum-free sets.

Fredricksen and Sweet gave \(S(6)\ge 536\) {[}FS00{]}, and Chen, Zhu and
Long gave \(S(6)\ge 537\) {[}CZL26{]}. The latter is the largest lower
bound in the sources that we read on 11 November 2026.

\textbf{Theorem 1.} \(S(6)\ge 538\).

\emph{Proof.} Let \(W\) be the word of 538 letters in Figure 1, and for
\(i=1,\dots,6\) let \(C_i\) be the set of positions \(n\) such that
letter \(n\) of \(W\) is \(i\). These six sets partition
\(\{1,\dots,538\}\), and their sizes are \(131, 85, 58, 67, 63, 134\).
There are \(269^2=72361\) pairs \((x,y)\) with \(1\le x\le y\) and
\(x+y\le 538\). For each pair, the letters at positions \(x\), \(y\) and
\(x+y\) are not all the same. So each \(C_i\) is sum-free. \(\square\)

Figure 1. The word \(W\). Each row gives the position of its first
letter.

\begin{verbatim}
  1  1214532151 4154552214 1213122344 2151413435 2315124412
 51  2541313532 2515141312 2334213142 5552231312 3452214421
101  5133531314 1552443221 3122432215 1514353321 4121614245
151  3215141545 5121412163 6214121614 5641516161 6266661413
201  1656465161 4166626616 1616266661 6661666666 6666666636
251  6161666666 6666666666 6666666666 6666666663 6666666166
301  6666161666 6662616666 6626661666 6564656131 4166662616
351  1615646541 6121412636 1214121564 5141512354 2416121412
401  3351415151 2234221312 2344255141 3135231512 4412254321
451  3132255524 1312433221 3141515223 5313145221 4421513253
501  4314151244 3221312141 2255451415 12354121
\end{verbatim}

Let \(R_k(3)\) be the least \(N\) such that every colouring of the edges
of the complete graph \(K_N\) with \(k\) colours has a monochromatic
triangle.

\textbf{Corollary 2.} \(R_6(3)\ge 540\).

\emph{Proof.} We use the classical construction behind
\(S(k)\le R_k(3)-2\) {[}FS00, (2){]}. Take the vertices
\(0,1,\dots,538\) of \(K_{539}\), and give the edge \(\{a,b\}\) the
colour \(i\) with \(|a-b|\in C_i\). A monochromatic triangle \(a<b<d\)
would put \(x=b-a\), \(y=d-b\) and \(x+y=d-a\) in one set \(C_i\). This
is not possible, because \(C_i\) is sum-free. So this colouring with six
colours has no monochromatic triangle, and neither has its restriction
to fewer vertices. So \(R_6(3)>539\). \(\square\)

Corollary 2 improves the bound \(R_6(3)\ge 539\) of {[}CZL26{]}.

\textbf{Corollary 3.} \(S(8)\ge 5382\), \(S(11)\ge 204588\) and
\(S(13)\ge 2045882\).

\emph{Proof.} We use the recurrences \(S(k+2)\ge 10\,S(k)+2\) {[}BBG26,
Section IV{]} and \(S(k+5)\ge 380\,S(k)+148\) {[}ACP22, inequality
(6){]}. With Theorem 1 they give
\[S(8)\ge 10\cdot 538+2=5382,\qquad S(11)\ge 380\cdot 538+148=204588,\]
\[S(13)\ge 10\cdot 204588+2=2045882. \qquad\square\]

These improve the bounds \(5372\), \(204208\) and \(2042082\) of
{[}CZL26{]}, which come from the same recurrences applied to
\(S(6)\ge 537\).

\section{2. Checks}\label{checks}

The following Python program holds the only computation that Theorem 1
needs. With \texttt{W} the word of Figure 1, as one string, it prints
\texttt{538\ {[}\textquotesingle{}1\textquotesingle{},\ \textquotesingle{}2\textquotesingle{},\ \textquotesingle{}3\textquotesingle{},\ \textquotesingle{}4\textquotesingle{},\ \textquotesingle{}5\textquotesingle{},\ \textquotesingle{}6\textquotesingle{}{]}\ {[}{]}}.
The last item is the list of pairs \((x,y)\) with letters at \(x\),
\(y\) and \(x+y\) all the same.

\begin{Shaded}
\begin{Highlighting}[]
\NormalTok{w }\OperatorTok{=} \StringTok{" "} \OperatorTok{+}\NormalTok{ W}
\NormalTok{bad }\OperatorTok{=}\NormalTok{ [(x, y) }\ControlFlowTok{for}\NormalTok{ x }\KeywordTok{in} \BuiltInTok{range}\NormalTok{(}\DecValTok{1}\NormalTok{, }\DecValTok{270}\NormalTok{) }\ControlFlowTok{for}\NormalTok{ y }\KeywordTok{in} \BuiltInTok{range}\NormalTok{(x, }\DecValTok{539} \OperatorTok{{-}}\NormalTok{ x)}
       \ControlFlowTok{if}\NormalTok{ w[x] }\OperatorTok{==}\NormalTok{ w[y] }\OperatorTok{==}\NormalTok{ w[x }\OperatorTok{+}\NormalTok{ y]]}
\BuiltInTok{print}\NormalTok{(}\BuiltInTok{len}\NormalTok{(W), }\BuiltInTok{sorted}\NormalTok{(}\BuiltInTok{set}\NormalTok{(W)), bad)}
\end{Highlighting}
\end{Shaded}

The existence of the partition is also checked in Lean 4, version
v4.35.0-rc3. The file \texttt{lean/SchurSix538.lean} imports Lean core
only and states, in the namespace \texttt{SchurWitness538}:

\begin{Shaded}
\begin{Highlighting}[]
\NormalTok{public theorem schur\_six\_538 :}
\NormalTok{    ∃ c : Nat → Fin 6, ∀ x y : Nat, 1 ≤ x → 1 ≤ y → x + y ≤ 538 →}
\NormalTok{      ¬ (c x = c y ∧ c y = c (x + y))}
\end{Highlighting}
\end{Shaded}

For the witness in the file, \(c(n)+1\) is letter \(n\) of \(W\) for
\(1\le n\le 538\) (theorem \texttt{colouring538\_val}, which is private
to the file). So the sets \(\{n\in\{1,\dots,538\}: c(n)=i\}\), for \(i\)
in \texttt{Fin\ 6}, are the sets \(C_{i+1}\) of Theorem 1. Lean does not
define \(S(k)\) or \(R_k(3)\): the step from the partition to
\(S(6)\ge 538\) is the definition of \(S(6)\), and Corollaries 2 and 3
are not formalized. The theorem uses the axioms \texttt{propext} and
\texttt{Quot.sound}. This file has no \texttt{sorry} and does not use
\texttt{native\_decide}. The file \texttt{lean/Challenge.lean} states
the same theorem with \texttt{sorry} on purpose, for the comparator of
the Lean package; the proof does not use it.

Corollary 3 rests on two preprints, but its three partitions are also
checked directly, and this check uses only the tables of those
preprints. The program \texttt{check\_recursive\_bounds.py} in the
repository applies {[}BBG26, Table II{]} to \(W\) (8 sets,
\(\{1,\dots,5382\}\)), {[}ACP22, Table 7{]} to \(W\) (11 sets,
\(\{1,\dots,204588\}\)), and {[}BBG26, Table II{]} to the partition with
11 sets (13 sets, \(\{1,\dots,2045882\}\)). It finds that each set is
sum-free.

\section{3. Data availability}\label{data-availability}

The word, the Lean package and the program
\texttt{check\_recursive\_bounds.py} are at
\url{https://github.com/mysticflounder/schur-538}, under the license CC
BY 4.0. The file \texttt{schur6-538.word} has the letters of \(W\) and
one final newline, which must be removed to get the output stated in
Section 2. Its SHA-256 is

\texttt{9e149c007b778655afd199685ebcc76260604a44f4c322d8aa9a1d96afb71d20},

the SHA-256 of \texttt{lean/SchurSix538.lean} is

\texttt{ef4fa1757ff30ed191f7eb1af24af288ab0a67726ee52b108355402789a63b81},

and the SHA-256 of \texttt{check\_recursive\_bounds.py} is

\texttt{aa7620b4dceb9788fa2a1babca998b01f2bcf266ac1fc9aaeda896397f25ff0c}.

The repository also holds a second such word,
\texttt{schur6-538-b.word}, which differs from \(W\) in 37 letters. The
Lean theorem does not use it.

\section{Use of AI tools}\label{use-of-ai-tools}

The word was found by a computer program whose development the author
directed. AI models (Anthropic Claude and OpenAI GPT models) wrote the
program. The computer checks, the Lean formalization and this note were
produced with the AI system Claude (Anthropic; model Claude Opus 5.5)
under the direction of the author. Separate agents of the same model
audited the word, the Lean theorem and the partitions of Corollary 3
with their own code. The author has checked the content and takes
responsibility for it.

\section{References}\label{references}

\begin{itemize}
\tightlist
\item
  {[}ACP22{]} R. Ageron, P. Casteras, T. Pellerin, Y. Portella, A.
  Rimmel, J. Tomasik, New lower bounds for Schur and weak Schur numbers,
  preprint, arXiv:2112.03175v2 (2022).
\item
  {[}BBG26{]} N. Bengone, A. Brouk, M. Grinsztajn, T. Helbert, B.
  Lugherini, A. Rimmel, J. Tomasik, Shifted S-templates and improved
  lower bounds for Schur numbers, preprint, arXiv:2607.15034v1 (2026).
\item
  {[}CZL26{]} Z. Chen, L. Zhu, F. Long, A new lower bound for the sixth
  Schur number: \(S(6)\ge 537\), preprint, arXiv:2610.31362v1 (2026).
\item
  {[}FS00{]} H. Fredricksen, M. M. Sweet, Symmetric sum-free partitions
  and lower bounds for Schur numbers, Electron. J. Combin. 7 (2000),
  \#R32. doi:10.37236/1510.
\end{itemize}

\end{document}
